\usepackage{fontspec}
\usepackage{polyglossia}
\setdefaultlanguage{russian}
\setotherlanguage{english}

% The manual is primarily built on Windows/MiKTeX.  Probe the standard
% Windows fonts first so MiKTeX does not try to synthesize missing Noto/
% Liberation fonts through makemf/hbf2gf.  The explicit Polyglossia
% Cyrillic font families are required even when the selected main font
% already contains Cyrillic glyphs.
\IfFontExistsTF{Times New Roman}{
  \setmainfont{Times New Roman}
  \setsansfont{Arial}
  \setmonofont{Consolas}[Scale=MatchLowercase]
  \newfontfamily\cyrillicfont{Times New Roman}[Script=Cyrillic]
  \newfontfamily\cyrillicfontsf{Arial}[Script=Cyrillic]
  \newfontfamily\cyrillicfonttt{Consolas}[Script=Cyrillic,Scale=MatchLowercase]
}{
  \IfFontExistsTF{Liberation Serif}{
    \setmainfont{Liberation Serif}
    \setsansfont{Liberation Sans}
    \setmonofont{Liberation Mono}[Scale=MatchLowercase]
    \newfontfamily\cyrillicfont{Liberation Serif}[Script=Cyrillic]
    \newfontfamily\cyrillicfontsf{Liberation Sans}[Script=Cyrillic]
    \newfontfamily\cyrillicfonttt{Liberation Mono}[Script=Cyrillic,Scale=MatchLowercase]
  }{
    \setmainfont{DejaVu Serif}
    \setsansfont{DejaVu Sans}
    \setmonofont{DejaVu Sans Mono}[Scale=MatchLowercase]
    \newfontfamily\cyrillicfont{DejaVu Serif}[Script=Cyrillic]
    \newfontfamily\cyrillicfontsf{DejaVu Sans}[Script=Cyrillic]
    \newfontfamily\cyrillicfonttt{DejaVu Sans Mono}[Script=Cyrillic,Scale=MatchLowercase]
  }
}

\usepackage[margin=2.2cm]{geometry}
\usepackage{microtype}
\usepackage{amsmath}
\usepackage{booktabs}
\usepackage{longtable}
\usepackage{array}
\usepackage{enumitem}
\usepackage{xcolor}
\usepackage{listings}
\usepackage{hyperref}
\usepackage{bookmark}
\usepackage{csquotes}

\hypersetup{
  unicode=true,
  colorlinks=true,
  linkcolor=blue,
  urlcolor=blue,
  citecolor=blue,
  pdftitle={Grandgousier — пользовательское руководство},
  pdfauthor={Проект Grandgousier}
}

\setlist{nosep}
\setlength{\parindent}{1.25em}
\setlength{\parskip}{0.35em}
\emergencystretch=2em

\lstdefinestyle{grandgousier}{
  basicstyle=\ttfamily\small,
  breaklines=true,
  columns=fullflexible,
  keepspaces=true,
  frame=single,
  showstringspaces=false
}
\lstset{style=grandgousier}

\newcommand{\program}{\texttt{Grandgousier}}
\newcommand{\cli}[1]{\texttt{#1}}
\newcommand{\stage}[1]{\textbf{Стадия #1}}
